\documentclass[10pt,a4paper]{amsart}
\usepackage[margin=19mm]{geometry}
\usepackage{amsmath,amssymb,mathtools}
\usepackage[scr=rsfs,cal=euler]{mathalfa}
\usepackage{xfrac}
\usepackage[unicode,hidelinks,bookmarks=false]{hyperref}
\newcommand{\ie}{i.e.\ }
\newcommand{\cf}{cf.\ }
\newcommand{\resp}{respectively\ }
\newcommand{\Subsection}{\subsection}
\providecommand{\nobreakdash}{\nobreak\hbox{-}}
\setlength{\emergencystretch}{2em}
\multlinegap0pt
\usepackage{enumerate}
\usepackage{xcolor}
\newtheorem{lem}{Lemma}
\newtheorem{defi}{Definition}
\def\P{\mathcal{P}}
\def\PB{\partial B_1}
\def\R {\mathbb{R}}
\def\SP{\mathcal{SP}}
\def\hem{\hspace{0.5em}}
\def\vem{\vspace{0.6em}}
\title{Concentration of cones in the Alt-Phillips problem}
\author{Ovidiu Savin, Hui Yu}
\date{Source: arXiv:2503.03626v2 (CC BY 4.0)}
\begin{document}
\maketitle

\noindent\emph{Source and licence.} Concentration of cones in the Alt-Phillips problem. Version arXiv:2503.03626v2 (CC BY 4.0), distributed by arXiv under the Creative Commons Attribution 4.0 International licence, http://creativecommons.org/licenses/by/4.0/, as recorded in the arXiv licence metadata for this exact version and shown on the abstract page https://arxiv.org/abs/2503.03626v2. Full licence notice: \texttt{licenses/arxiv-2503.03626-CC-BY-4.0.txt}.\par\smallskip

\noindent\textit{Editorial scope.} Counted unit: the article's lem LemMagicLemma (source0.tex lines 189--200). The article's complete original proof of it is delivered with it (source0.tex lines 529--598, one \texttt{proof} environment of 70 lines, which the article places away from the statement). No mathematics is written, repaired or re-derived: the statement and the proof are the article's own lines, in the article's own notation, and no retained line is substituted or re-flowed.  Retained uncounted as the source-local context the proof needs: lines 137--140; lines 162--165; lines 362--376; lines 404--417; lines 456--459; lines 467--470; lines 477--486.  Nothing was omitted from any retained span, so the package carries no omission record.  No retained line contains a citation, so the wrapper carries no bibliography and no bibliography entry.  No retained line contains a figure environment, an image-inclusion command or drawing code, so no image is delivered and the package declares no asset dependency.  Editorial formatting only, in addition: an AMS \texttt{amsart} preamble that restates the article's own declarations and macros used by the retained lines, namely the environments \texttt{lem}, \texttt{defi} on separate counters, together with the standard packages those lines need.  The retained environments print in this document as Lemma 1, Definition 1 in order of appearance, whereas the article prints the same items with the numbering of its own full document.  The complete native source of the article as distributed in the arXiv e-print for this exact version is bundled unchanged as \texttt{sources/174276/source0.tex}.
\begin{equation}
\label{EqnParabolaSolutions}
\frac12x\cdot Ax, \text{ where  }A\ge0 \text{ and }\mathrm{trace}(A)=1.
\end{equation}

\begin{equation}
\label{EqnkCone}
P_k(x):=\frac{1}{2k}\sum_{1\le j\le k}x_j^2 \hem\text{ for }k=1,2,\dots,d.
\end{equation} 

\begin{defi}
\label{DefObPCones}
The collection of half-space solutions for the classical obstacle problem is denoted by $\mathcal{HS}$, namely,
$$
\mathcal{HS}:=\{\frac12[(x\cdot e)_+]^2: \hem e\in\R^d, \hem |e|=1\}.
$$
The collection of parabola solutions is denoted by $\mathcal{P}$, namely,
$$
\mathcal{P}:=\{\frac12x\cdot Ax:\hem A\ge0,\hem \mathrm{trace}(A)=1\}.
$$
The collection of parabola solutions with symmetry is denoted by $\mathcal{SP}$, namely,
$$
\mathcal{SP}:=\{p\in\mathcal{P}: \hem \text{up to a rotation,}\hem p=P_k \text{ for some }k\in\{1,2,\dots,d\} \text{ from }\eqref{EqnkCone}\}.
$$
\end{defi}

\begin{lem}
\label{LemIntegralInDifferentDimensions}
For $d\ge2,$ suppose that $p\in\P$ is invariant in the $e_d$-direction, that is, 
$$
p(x_1,x_2,\dots,x_d)=q(x_1,x_2,\dots,x_{d-1}),
$$
for some $q\in\P$ in $\R^{d-1}$.

Then 
$$
\int_{\partial B_1}\frac{|\nabla p|^2}{p}(p-\frac{1}{2d}|x|^2)d\mathcal{H}^{d-1}=\alpha_d\int_{\partial B_1\cap\{x_d=0\}}\frac{|\nabla q|^2}{q}\cdot(q-\frac{1}{2(d-1)}|(x_1,x_2,\dots x_{d-1})|^2)d\mathcal{H}^{d-2}
$$
for a positive dimensional constant $\alpha_d$.
\end{lem} 

\begin{equation}
\label{EqnInterpolation}
p_t(x):=tp(x)+(1-t)P_d(x)\hem \text{ for }t\ge0.
\end{equation} 

\begin{equation}
\label{EqnQInterpolation}
Q(t):=\int_{\partial B_1}\frac{|\nabla p_t|^2}{p_t}(p_t-\frac{1}{2d}|x|^2) \hem\text{ as long as }p_t\ge0. 
\end{equation} 

\begin{lem}
\label{LemMeaningfulQ}
For $p\in\P$, let $p_t$ and $Q$ be as in \eqref{EqnInterpolation} and \eqref{EqnQInterpolation}. Then we have
$$
Q(t)=t^2\cdot[\int_{\partial B_1}|\nabla_\tau p|^2+4\int_{\partial B_1}(p-P_d)^2-\frac{1}{2d}\int_{\partial B_1}\frac{|\nabla_\tau p|^2}{p_t}]
$$
as long as $p_t\ge0.$ 

Here $\nabla_\tau$ denotes the gradient along directions tangential to $\PB$.
\end{lem} 

\begin{lem}
\label{LemMagicLemma}
Suppose that $p$ is a parabola solution to the classical obstacle problem in $\R^d$ from \eqref{EqnParabolaSolutions}, then
\begin{equation*}
\int_{\partial B_1}\frac{|\nabla p|^2}{p}(p-\frac{1}{2d}|x|^2)\ge0.
\end{equation*} 
The equality  holds if and only if, up to a rotation, 
$$
p=P_k
$$
 for some $k=1,2,\dots,d$ from \eqref{EqnkCone}.
\end{lem} 

\begin{proof}[Proof of Lemma \ref{LemMagicLemma}]
We argue by an induction on the dimension $d$. 

In $\R^1$, there is only one parabola solution, namely, $x^2/2$. The conclusion follows. 

Suppose that we have established Lemma \ref{LemMagicLemma} in dimensions $1,2,\dots, d-1$, below we prove the result in $\R^d$. 

\vem

For $p_t$ from \eqref{EqnInterpolation}, define 
\begin{equation}
\label{EqnBT}
\bar{t}:=\sup\{t:\hem p_t\ge0 \text{ on }\PB\}.
\end{equation}  
Since  $p\in\P$ (see Definition \ref{DefObPCones}), we have $$\bar{t}\ge1.$$ 

Moreover, by definition of $\bar{t}$, we have
$$
p_{\bar t}\in\P \hem \text{ with }\hem p_{\bar t}(e)=0 
$$
at some $e\in\PB$. This implies that $p_{\bar t}$ is invariant along the $e$-direction. Our induction hypothesis and   Lemma \ref{LemIntegralInDifferentDimensions} give, for the quantity $Q$ in \eqref{EqnQInterpolation},
\begin{equation}
\label{EqnInductionHypothesis}
Q(\bar t)\ge 0, \hem\text{ and equality holds only when }p_{\bar t}\in\SP.
\end{equation} 
Recall the space of symmetric cones $\SP$ from Definition \ref{DefObPCones}. 

\vem
Based on the value of $\bar t$ in \eqref{EqnBT}, we consider two cases.

If $\bar{t}=1$, then $p=p_{\bar t}$. The desired conclusion follows by \eqref{EqnInductionHypothesis}.  

It remains to consider the case when $\bar{t}>1$.

In this case, define
\begin{equation}
\label{Eqnq}
q(t):=Q(t)/t^2=\int_{\partial B_1}|\nabla_\tau p|^2+4\int_{\partial B_1}(p-P_d)^2-\frac{1}{2d}\int_{\partial B_1}\frac{|\nabla_\tau p|^2}{p_t},
\end{equation}
where we used Lemma \ref{LemMeaningfulQ}. 

With \eqref{EqnInterpolation}, we have $p_0=\frac{1}{2d}$ on $\PB$. This implies
\begin{equation}
\label{Eqnq0}
q(0)=4\int_{\partial B_1}(p-P_d)^2\ge0.
\end{equation}
At $\bar t>1$, we have
\begin{equation}
\label{Eqnqbt}
q(\bar t)=Q(\bar t)/\bar{t}^2\ge0
\end{equation}
by \eqref{EqnInductionHypothesis}.

Moreover, direct computation with \eqref{EqnInterpolation} and \eqref{Eqnq} gives, for $t\in(0,\bar t)$,
$$
q''(t)=-\frac{1}{d}\int_{\PB}\frac{|\nabla_\tau p|^2}{p_t^3}(p-P_d)^2\le 0.
$$
Together with \eqref{Eqnq0} and \eqref{Eqnqbt}, this implies 
$$q(1)\ge0.$$
With \eqref{EqnQInterpolation} and \eqref{Eqnq}, this gives the desired inequality 
\begin{equation}
\label{EqnDesiredInequality}
\int_{\PB}\frac{|\nabla p|^2}{p}(p-\frac{1}{2d}|x|^2)=Q(1)=q(1)\ge0.
\end{equation}

The equality case in \eqref{EqnDesiredInequality} forces the equality in \eqref{Eqnq0}, this implies
$$
p=P_d=\frac{1}{2d}|x|^2.
$$
\end{proof} 

\end{document}
